\lyseis{Α4.2}{Κανονικότητες της μορφής α·ν}
\sol{1} Οι επόμενοι όροι είναι $45$ και $54$. Γενικός όρος ο $9\nu$ (β): ο $1$ος όρος είναι $9\cdot1=9$ και κάθε όρος αυξάνεται κατά $9$. Για τη σειρά $4,8,12,16$: ο τύπος $4\nu$ (β) δίνει $4,8,12,16$. Οι άλλοι δεν κάνουν: $4+\nu$ δίνει $5$ για $\nu=1$, $4+4\nu$ δίνει $8$, και $8\nu-4$ δίνει $4$, $12$.
\sol{2} α) $8,\ 16,\ 24,\ 32,\ 40$. β) $11,\ 22,\ 33,\ 44,\ 55$. γ) $-3,\ -6,\ -9,\ -12,\ -15$. δ) $0{,}5,\ 1,\ 1{,}5,\ 2,\ 2{,}5$.
\sol{3} α) $30,\ 36$ και γενικός όρος $6\nu$. β) $-10,\ -12$ και $-2\nu$. γ) $1{,}25,\ 1{,}5$ και $0{,}25\nu$ (ο πρώτος όρος είναι $0{,}25$ και το βήμα $0{,}25$).
\sol{4} $9\cdot10=90$, $9\cdot25=225$, $9\cdot100=900$.
\sol{5} $6\nu=72$, $\nu=72:6=12$ (ο $12$ος). $6\nu=240$, $\nu=40$ (ο $40$ός). $6\nu=1\,500$, $\nu=250$ (ο $250$ός).
\sol{6} α) Ναι: $91:7=13$, είναι ο $13$ος όρος. β) Όχι: $150:7=21{,}4\ldots$ δεν είναι φυσικός. γ) Ναι: $700:7=100$, ο $100$ός όρος. δ) Όχι: $5:7<1$, δεν είναι φυσικός.
\sol{7} α) Λάθος: είναι $5\nu$ (το $5+\nu$ δίνει $6$ για $\nu=1$). β) Σωστό: $3\cdot4=12$. γ) Σωστό. δ) Σωστό. ε) Λάθος: το $\nu$ είναι θέση, φυσικός αριθμός. στ) Σωστό: $\alpha\cdot1=\alpha$. ζ) Σωστό: $2\cdot9=18$.
\sol{8} α) $3\cdot7=21$ οδοντογλυφίδες. β) $3\nu=120$, $\nu=40$: πλευρά $40$. γ) $3\nu=100$, $\nu=100:3=33{,}3\ldots$, που δεν είναι φυσικός: όχι.
\sol{9} Για $3,4,5,6$ πλευρές οι ακμές είναι $6,8,10,12$. Για $8$-γωνο $16$, για $9$-γωνο $18$ και για $20$-γωνο $40$. Γενικός όρος: $2\nu$ (το διπλάσιο των πλευρών της βάσης).
\sol{10} $\nu$ παγωτά κοστίζουν $1{,}5\nu$ €. Τα $12$ κοστίζουν $1{,}5\cdot12=18$ €. Με $27$ €: $1{,}5\nu=27$, $\nu=27:1{,}5=18$ παγωτά.
\sol{11} Μετά από $12$ δευτερόλεπτα: $\frac{3}{4}\cdot12=9$ μέτρα. Για τα $15$ μέτρα: $\frac{3}{4}\nu=15$, $\nu=15\cdot\frac{4}{3}=20$ δευτερόλεπτα.
\sol{12} Η θερμοκρασία μετά από $\nu$ ώρες είναι $-1{,}5\nu$. Μετά από $8$ ώρες: $-1{,}5\cdot8=-12$°. Για τους $-27$°: $-1{,}5\nu=-27$, $\nu=(-27):(-1{,}5)=18$: μετά από $18$ ώρες.
\sol{13} Για $\nu=1$ ο τύπος του Στέλιου δίνει $6+6\cdot1=12$, ενώ ο $1$ος όρος είναι $6$. Ο σωστός τύπος είναι $6\nu$.
\sol{14} $\alpha\cdot4=10$, άρα $\alpha=10:4=2{,}5$. Οι δύο πρώτοι όροι είναι $2{,}5$ και $5$. Ο $50$ός είναι $2{,}5\cdot50=125$.
\sol{15} Πρέπει $\alpha\nu=30$ με $\nu$ φυσικό, δηλαδή ο $\alpha$ να είναι διαιρέτης του $30$: $\alpha=1,\,2,\,3,\,5,\,6,\,10,\,15,\,30$.
\sol{16} Η σειρά της Ζωής έχει γενικό όρο $3\nu$, γιατί ο πρώτος όρος ($3$) ισούται με το βήμα. Η σειρά του Ηλία έχει βήμα $3$, αλλά ο πρώτος όρος είναι $4$: τα $4,\,7,\,10,\dots$ είναι $3\nu+1$, δεν έχουν τη μορφή $\alpha\nu$ ($3\cdot1=3\neq4$).
\sol{17} Ελεύθερη απάντηση. Π.χ. $12\nu$: αυγά σε $\nu$ δωδεκάδες, μήνες σε $\nu$ έτη. $0{,}5\nu$: λίτρα νερού για $\nu$ άτομα με μισό λίτρο το καθένα, ή ευρώ για $\nu$ καραμέλες των $0{,}5$ €.
\sol{18} Οι όροι είναι $2{,}5,\ 5,\ 7{,}5,\ 10,\ 12{,}5,\ 15,\ 17{,}5,\ 20$. Τα σημεία $(1;2{,}5)$, $(2;5)$, $\dots$ βρίσκονται σε μια ευθεία γραμμή που περνά από την αρχή των αξόνων.
